% paper_vlsid.tex — VLSID submission draft (deadline Sep 10)
% Story: unmasked NIST ASCON-128 hardware core on CW305:
%   (1) thorough profiled DL-SCA fails at the VCCINT measurement floor,
%   (2) clock-glitch FI produces real, localized finalization faults,
%   (3) key recovery: [MONDAY OUTCOME SLOT — see KEY RESULT macros]
% Structured so Monday's board results drop into the marked slots only.
% Class: local texlive has no IEEEtran; article-twocolumn compiles here
% (same as paper.tex). Swap to the VLSID template at submission.
\documentclass[a4paper,11pt,twocolumn]{article}

\usepackage[margin=2cm]{geometry}
\usepackage{amsmath,amssymb}
\usepackage{booktabs}
\usepackage{graphicx}
\usepackage{hyperref}
\usepackage{cite}

% ── KEY RESULT MACROS — filled after Monday ──────────────────────────
% Outcome ladder (choose ONE at submission):
%   \blackboxtrue  : Track B/C full key on the unmodified core (headline)
%   \demoonlytrue  : instrument-guided study + disclosure demo + sim solvers
%   \negativeonlytrue : negative + methodology paper (floor)
\newif\blackbox
\newif\demoonly
\newif\negativeonly
\negativeonlytrue            % placeholder until Monday decides

\title{Clock-Glitch Fault Injection and Profiled Side-Channel Analysis of
the NIST ASCON-128 Lightweight Cryptography Hardware Reference:
A Measurement-Floor Study with Verified Fault Localization}

\author{Prajwal R\\nPES University, ECE, Bengaluru, India\\nruyketsu.prajwal@gmail.com}}

\begin{document}
\maketitle

\begin{abstract}
% [MONDAY: 2-3 sentences keyed to outcome ladder]
We study the security of the unprotected NIST ASCON-128 hardware reference
CryptoCore on a CW305 (Artix-7) FPGA target under two attack families:
profiled deep-learning power side-channel analysis and clock-glitch fault
injection. Power analysis is defeated by the measurement channel rather
than the algorithm: total-core-current (VCCINT) probing of a 100k-gate FPGA
places the 3k-LUT cipher's leakage at $-19$\,dB SNR, and we characterize
the floor exhaustively across gains, clock rates, M-averaging, and
architectures (CNN/MLP/transformer), with a single S-box column's two key
bits recoverable only at $M{=}64$ trace averaging ($+18$\,dB). Clock
glitching by contrast injects real, reproducible, single-bit faults into
the finalization permutation; we localize them to specific permutation
stages from public outputs alone (tag-differential signatures), and
\ifblackbox recover the full 128-bit key by differential fault analysis
in $N{=}$[FILL] queries, verified against the public-output oracle.
\fi
\ifdemoonly
demonstrate single-query full-key disclosure under an explicit
state-observability threat model on an instrumented build, with the
black-box differential-fault-analysis pipeline validated end-to-end in
simulation (15 faults/column $\rightarrow$ full key).
\fi
\ifnegativeonly
characterize the fault-placement barrier that separates the achieved
faults from differential-fault-analyzable ones.
\fi
We document three methodology traps that produce confident false
positives in this pipeline and release all capture/attack tooling.
\end{abstract}

\section{Introduction}
% Hook: NIST LWC winner deployed in constrained HW; its official hardware
% reference must withstand both SCA and FI. We test both on real silicon.
% Contributions (numbered):
%  C1: exhaustive PA-negative at the VCCINT measurement floor (+ the one
%      column that DOES converge at M=64 — quantifies exactly what the
%      floor costs the attacker).
%  C2: first clock-glitch faults inside ASCON's finalization on this
%      target, localized to permutation stage from public outputs only.
%  C3: [MONDAY SLOT] full-key result: black-box DFA / observability demo /
%      barrier characterization (per outcome ladder).
%  C4: three leakage-gate false-positive traps (prior-only CNN, edge gate,
%      nonce-aliasing) + released toolchain.

\section{Background and Target}
% ASCON-128 structure: init p^12 (state = IV||K||N), absorb, finalization
% p^6 with key whitening; tag = S3^k0 || S4^k1.
% Target: ascon-hardware-sca unprotected CryptoCore (VHDL, CCW=32),
% CW305-100t, CC-Husky scope, 40 MS/s, tio4 trigger, 10 MHz crypto clock.
% Register map + adapter. [keep the honesty note: state observability port
%  is a local addition; official design exposes no such port.]

\section{Power Side-Channel Analysis: Characterizing the Floor}
% 3.1 setup: random-key captures, M-averaging, gains, clock rates
% 3.2 profiling: CNN/MLP/joint/transformer; floor-gated columns
% 3.3 results table: SNR per config; top-1 vs chance; M-sweep
%     (+18 dB at M=64; single column converges, 2 bits, posterior 0.999)
% 3.4 KADD full-key intermediate: 2x-chance leakage, +3.1 nats/trace for
%     the true key, unfactorable (12 rounds diffusion) — structural wall
% 3.5 verdict: measurement channel, not algorithm; what would change it
%     (EM near-field, dedicated core supply)

\section{Clock-Glitch Fault Injection}
% 4.1 injection path: Husky glitch module -> 20-pin -> tio_clkin; setup
%     traps (phase-step units, arm-before-go, DCM-relock artifacts — the
%     settled-scan methodology)
% 4.2 fault evidence: deterministic single-bit finalization faults,
%     ct-intact; eo=31 = finalization-entry (key-injection stage) fault,
%     identified by exhaustive model matching from the public faulty tag
% 4.3 tag-differential signature as a stage oracle: 62-bit = early,
%     1-bit = output latch, 5-45 bit = round-11 S-box input (DFA band)
% 4.4 [MONDAY SLOT: cycle-map-guided hunt + nonce/voltage sweep results]

\section{Key Recovery}
% [MONDAY SLOT — one of:]
%  5A black-box: round-11 faults -> masked-DDT DFA -> 64 columns -> key,
%     verified vs public oracle; query/fault counts; OR
%     PA+assembly: >=52/64 confident columns -> bounded 4^k brute force
%     -> oracle-verified key (selftest 30/30, budget table)
%  5B demo: glitch_only hang -> frozen state -> round inversion -> key,
%     self-verified (IV||nonce), one faulted query; explicit
%     state-observability threat model; instrument limitation stated
%  5C barrier: why the round-11 band stayed empty; sim-validated solvers
%     (DFA 15 f/col; assembly margins) as the methodology contribution

\section{Methodology Traps in Profiled SCA}
% T1: edge gate is prior-driven; T2: prior-only CNN + separating-nonce
% picker = confident wrong convergence (hyp 3 vs true 0 case); T3:
% fixed-key labels alias public nonce bits. All three with the concrete
% numbers from our logs; floor-gating as the fix.

\section{Related Work}
% ascon-hardware-sca TVLA baselines; ASCON DFA (bit-flip, eprint 2023/1923,
% Nakamura et al.); SIFA/SHFA line; CW305 glitching tutorials; debug-port
% attack class for the observability demo framing.

\section{Conclusion}
% [MONDAY: one paragraph per outcome ladder]

\section*{Artifact Availability}
% capture/training/FI scripts, solvers, selftests, raw fault npz
% (repo on request / anonymized for review)

\end{document}
